% ==============================================================================
% Préambule LaTeX Beamer - ANF ArchiDrone 2026 (Drone&Cap' / 2RM)
% Auteur : Quentin Brateau
% Thématique : Drones Marins & Sous-Marins pour la Recherche Scientifique
% ==============================================================================

\usepackage[utf8]{inputenc}
\usepackage[T1]{fontenc}
\usepackage[french]{babel}

\usepackage{amsmath, amssymb, amsfonts}
\usepackage{booktabs}
\usepackage{tabularx}
\usepackage{colortbl}
\usepackage{graphicx}
\graphicspath{{../../imgs/}{../imgs/}{imgs/}}
\usepackage{tikz}
\usetikzlibrary{arrows.meta, positioning, calc, shapes.geometric, backgrounds, decorations.pathmorphing, shadows}

\usepackage{lmodern}
\usepackage{listings}
\usepackage{microtype}
\usepackage{hyperref}

% ------------------------------------------------------------------------------
% Palette de couleurs océanographiques modernes
% ------------------------------------------------------------------------------
\definecolor{OceanDeep}{HTML}{0A2240}      % Bleu profond abyssal
\definecolor{OceanMedium}{HTML}{0E4D7A}    % Bleu pélagique
\definecolor{OceanLight}{HTML}{1976D2}     % Bleu clair / surface
\definecolor{OceanCyan}{HTML}{00A8B5}      % Cyan lagon / accents
\definecolor{OceanCoral}{HTML}{FF6B6B}     % Corail (alertes, fail-safes)
\definecolor{OceanGreen}{HTML}{10B981}     % Vert posidonie (validations)
\definecolor{OceanGray}{HTML}{F0F4F8}      % Fond clair blocs
\definecolor{DarkSlate}{HTML}{1E293B}      % Texte sombre
\hypersetup{colorlinks=true, urlcolor=OceanLight, linkcolor=OceanMedium}

% ------------------------------------------------------------------------------
% Configuration du thème Beamer (Metropolis / mtheme)
% ------------------------------------------------------------------------------
\usetheme[
  progressbar=frametitle,
  block=fill,
  numbering=fraction
]{metropolis}


\setbeamercolor{palette primary}{bg=OceanDeep, fg=white}
\setbeamercolor{palette secondary}{bg=OceanMedium, fg=white}
\setbeamercolor{palette tertiary}{bg=OceanCyan, fg=white}
\setbeamercolor{structure}{fg=OceanMedium}

% Couleurs du texte et titres Metropolis
\setbeamercolor{normal text}{fg=DarkSlate, bg=white}
\setbeamercolor{title}{fg=OceanDeep}
\setbeamercolor{subtitle}{fg=OceanMedium}
\setbeamercolor{author}{fg=DarkSlate}
\setbeamercolor{date}{fg=DarkSlate}
\setbeamercolor{institute}{fg=DarkSlate}
\setbeamercolor{frametitle}{bg=OceanDeep, fg=white}
\setbeamercolor{section title}{fg=OceanDeep}

% Barres de progression et séparateurs Metropolis
\setbeamercolor{progress bar}{fg=OceanCyan, bg=OceanDeep!20}
\setbeamercolor{progress bar in head/foot}{fg=OceanCyan, bg=OceanDeep!20}
\setbeamercolor{progress bar in section page}{fg=OceanCyan, bg=OceanDeep!20}
\setbeamercolor{title separator}{fg=OceanCyan}

% Alertes et exemples
\setbeamercolor{alerted text}{fg=OceanCoral}
\setbeamercolor{example text}{fg=OceanGreen}

% Blocs personnalisés
\setbeamercolor{block title}{bg=OceanMedium, fg=white}
\setbeamercolor{block body}{bg=OceanGray, fg=DarkSlate}
\setbeamercolor{block title example}{bg=OceanGreen, fg=white}
\setbeamercolor{block body example}{bg=OceanGray, fg=DarkSlate}
\setbeamercolor{block title alerted}{bg=OceanCoral, fg=white}
\setbeamercolor{block body alerted}{bg=OceanCoral!10, fg=DarkSlate}
\setbeamerfont{block body}{size=\footnotesize}

% Table des matières (TOC)
\setbeamercolor{section in toc}{fg=OceanDeep}
\setbeamercolor{section in toc shaded}{fg=DarkSlate!30}

% Pied de page natif Metropolis
\setbeamertemplate{frame footer}{ANF ArchiDrone 2026 (Drone\&Cap' / 2RM)}

% Slide de transition à chaque nouvelle section (Style Standout)
\AtBeginSection[]{
  \begin{frame}[standout]
    \centering
    {\hypersetup{linkcolor=white}\LARGE\bfseries\textcolor{white}{\secname}\par}
    \vspace{-1.2em}
    {\color{OceanCyan}\rule{0.5\textwidth}{1.5pt}\par}
  \end{frame}
}

% ------------------------------------------------------------------------------
% Raccourcis et macros utiles
% ------------------------------------------------------------------------------
\newcommand{\tech}[1]{\textbf{\textcolor{OceanMedium}{#1}}}
\newcommand{\alertw}[1]{\textbf{\textcolor{OceanCoral}{#1}}}
\newcommand{\succw}[1]{\textbf{\textcolor{OceanGreen}{#1}}}
\newcommand{\code}[1]{\texttt{\small\textcolor{OceanDeep}{#1}}}

% Adaptation de \tech dans les titres de blocs pour conserver une parfaite lisibilité
\newcommand{\techinblock}[1]{\textbf{\textcolor{OceanCyan}{#1}}}
\newcommand{\techinblockwhite}[1]{\textbf{\textcolor{white}{#1}}}
\addtobeamertemplate{block begin}{\let\tech\techinblock}{}
\addtobeamertemplate{block alerted begin}{\let\tech\techinblockwhite}{}
\addtobeamertemplate{block example begin}{\let\tech\techinblockwhite}{}

% Ajustement fin des marges et du pied de page pour un rendu équilibré sans débordement
\makeatletter
\setlength{\metropolis@frametitle@padding}{1.5ex}
\setbeamertemplate{footline}{%
  \ifbool{metropolis@standout}{}{%
    \begin{beamercolorbox}[wd=\textwidth, sep=1.2ex]{footline}%
      \usebeamerfont{page number in head/foot}%
      \usebeamertemplate*{frame footer}%
      \hfill%
      \usebeamertemplate*{frame numbering}%
    \end{beamercolorbox}%
  }%
}
\makeatother

\newcommand{\ROV}{\textbf{ROV} (\textit{Remotely Operated Vehicle})}
\newcommand{\AUV}{\textbf{AUV} (\textit{Autonomous Underwater Vehicle})}
\newcommand{\ASV}{\textbf{ASV} (\textit{Autonomous Surface Vehicle})}
\newcommand{\USBL}{\textbf{USBL} (\textit{Ultra-Short BaseLine})}
\newcommand{\DVL}{\textbf{DVL} (\textit{Doppler Velocity Log})}
\newcommand{\CTD}{\textbf{CTD} (\textit{Conductivity-Temperature-Depth})}

